\section{Virtualização da Comunicação Entre Processadores}

A virtualização de interrupções de software na arquitetura RISC-V, um dos
mecanismos utilizados para a comunicação entre processadores, baseia-se na
intermediação das interrupções pelo hypervisor. Em uma implementação que não
dispõe da extensão H, o acesso ao mecanismo de interrupção de software
entre \textit{harts} é realizado em modo \textit{Machine}, por meio do
registrador \gls{msip}. Dessa forma, o \gls{so} convidado, executando em modo
\textit{Supervisor}, não possui acesso direto a esse registrador e,
consequentemente, não pode gerar diretamente uma interrupção de software para
outro processador físico. A geração e a entrega das interrupções são, portanto,
intermediadas pelo hypervisor, que as associa às respectivas CPUs virtuais.

Nessa implementação, o \gls{so} convidado solicita o envio de uma interrupção
por meio da interface de serviços definida pelo \textit{OpenSBI}. Ao interceptar
essa requisição, o hypervisor a direciona para a CPU virtual responsável pelo
processamento da chamada. A CPU virtual identifica as CPUs virtuais de destino
e encaminha a solicitação de interrupção. Quando a CPU virtual de destino está
em execução, o hypervisor utiliza uma interrupção de software física para
sinalizar o \textit{hart} correspondente. Após o processamento dessa
interrupção, a CPU virtual de destino registra a interrupção como pendente e a
entrega ao \gls{so}. Caso a CPU virtual de destino não esteja ativa, a interrupção é
apenas marcada como pendente. Nesse caso, sua entrega é realizada quando a CPU
virtual retomar a execução e seu estado for novamente atualizado pelo
hypervisor.

\begin{figure}[H]
\centering
\includegraphics[width=0.6\textwidth]{images/SoftwareInterruptVirtualization.pdf}
\caption{Diagrama de Envio de uma \gls{ipi} Virtual na RISC-V}\label{fig:virtual-ipi-riscv}
\end{figure}
